% Chapter 2: Background and Related Work
% Included via % Chapter 2: Background and Related Work
% Included via % Chapter 2: Background and Related Work
% Included via % Chapter 2: Background and Related Work
% Included via \include{chapter2} from report.tex

\chapter{Background and Related Work}
\label{ch:background}

This chapter establishes the technical foundations on which this thesis is built. Section~\ref{sec:esc} introduces the environmental sound classification task and the ESC-50 benchmark. Section~\ref{sec:snn-background} formalises the spiking neural network framework, covering the \gls{lif} neuron model, surrogate gradient training, and the seven spike encoding methods evaluated in this work. Section~\ref{sec:neuromorphic-hardware} describes the SpiNNaker neuromorphic platform and situates it among contemporary alternatives. Section~\ref{sec:energy-background} examines the energy efficiency argument for \glspl{snn}, including the conditions under which spiking architectures genuinely outperform conventional networks. Section~\ref{sec:related-audio-snn} surveys prior work on \gls{snn}-based audio classification and identifies the gap this thesis addresses. Finally, Section~\ref{sec:adversarial-cl} reviews findings on adversarial robustness and continual learning in \glspl{snn}, motivating two of the advanced experiments in this report.

%=================================================================
\section{Environmental Sound Classification}
\label{sec:esc}
%=================================================================

Environmental sound classification (ESC) is a core task within computational auditory scene analysis, requiring a system to assign a categorical label to a short audio recording of a real-world acoustic event~\citep{barchiesi2015}. Applications include noise pollution monitoring in smart cities~\citep{bello2019}, automated wildlife census~\citep{stowell2019}, assistive technology for the hearing-impaired, and anomaly detection in industrial settings. Unlike speech recognition, ESC must handle the full diversity of non-linguistic acoustic events---from a dog barking to a chainsaw operating---without the benefit of structured phonological constraints.

\paragraph{The ESC-50 dataset.}
\gls{esc}~\citep{piczak2015} is the standard benchmark for this task. It comprises 2{,}000 five-second recordings drawn from the Freesound repository, distributed uniformly across 50 classes (40 clips per class). The classes span five broad super-categories: animals (e.g.\ dog, rooster, crow), natural soundscapes (e.g.\ rain, sea waves, thunderstorm), human non-speech sounds (e.g.\ clapping, coughing, laughing), domestic/indoor sounds (e.g.\ clock alarm, door knock, vacuum cleaner), and urban/outdoor sounds (e.g.\ helicopter, siren, car horn). Crucially, the dataset defines five predefined folds for cross-validation, ensuring that results are directly comparable across studies. Human classification performance on ESC-50 is 81.3\%~\citep{piczak2015}, reflecting the genuine perceptual difficulty of some distinctions (e.g.\ sea waves versus rain, insects versus crackling fire).

The current state of the art is 99.1\%, achieved by OmniVec2~\citep{srivastava2024} using a multi-modal Transformer pretrained on diverse data. Among audio-only methods, Pretrained Audio Neural Networks~(\gls{panns}) with CNN14 fine-tuning achieve 94.7\%~\citep{kong2020}, and Audio Spectrogram Transformers reach 95.6\%~\citep{gong2021}. These results rely on large-scale pretraining on AudioSet~\citep{gemmeke2017}, a dataset of approximately 2 million 10-second clips spanning 527 event categories. Without external pretraining, the best reported \gls{ann} performance on ESC-50 is substantially lower, motivating our choice to evaluate scratch-trained models as the primary comparison.

\paragraph{Input representation.}
The dominant input representation for ESC is the log-mel spectrogram, computed by applying the short-time Fourier transform to the audio waveform, mapping the resulting frequency bins through a bank of triangular mel-scale filters that approximate the frequency resolution of the human cochlea, and converting power to a logarithmic decibel scale. In this work, we extract spectrograms with 64 mel bands at a sampling rate of 22{,}050\,Hz and a hop length of 512 samples, yielding a representation of shape $(1, 64, 216)$ per clip---one channel, 64 frequency bins, and 216 time frames. Values are normalised to $[0, 1]$ per sample, a requirement for the spike encoding methods described in Section~\ref{subsec:encoding}.

\paragraph{Broader context.}
Other datasets complement ESC-50. UrbanSound8K~\citep{salamon2014} provides 8{,}732 clips across 10 urban sound classes; it is smaller in vocabulary and widely considered less challenging. AudioSet~\citep{gemmeke2017} provides massive scale (527 classes, 2 million clips) but uses weak labels and variable-quality recordings, making it unsuitable as a controlled evaluation benchmark. The Spiking Heidelberg Digits (SHD) and Spiking Speech Commands datasets~\citep{cramer2020} are natively spiking but address spoken-digit recognition rather than environmental sounds. No standard spiking benchmark exists for ESC, underscoring the need for this work.


%=================================================================
\section{Spiking Neural Networks}
\label{sec:snn-background}
%=================================================================

\Glspl{snn} are neural networks in which information is communicated through discrete, binary events called spikes. Unlike conventional \glspl{ann}, where neurons transmit continuous-valued activations, \gls{snn} neurons accumulate incoming signals over time in a membrane potential and emit a spike only when that potential exceeds a threshold. This event-driven, temporal computation paradigm is inspired by biological neural circuits and offers potential advantages in energy efficiency and temporal processing~\citep{maass1997,gerstner2002}.

%-------------------------------
\subsection{The Leaky Integrate-and-Fire Neuron}
\label{subsec:lif}
%-------------------------------

The \gls{lif} neuron is the most widely used model in modern \gls{snn} research, balancing biological plausibility with computational tractability~\citep{gerstner2002}. Its continuous-time dynamics describe an RC circuit with passive membrane leakage:

\begin{definition}[Leaky Integrate-and-Fire Neuron]
\label{def:lif}
The membrane potential $U(t)$ of a \gls{lif} neuron evolves according to the differential equation
\begin{equation}
    \tau_m \frac{dU(t)}{dt} = -U(t) + R \, I(t),
    \label{eq:lif-continuous}
\end{equation}
where $\tau_m$ is the membrane time constant, $R$ is the membrane resistance, and $I(t)$ is the input current. When $U(t)$ reaches the firing threshold $\vartheta$, the neuron emits a spike and the membrane potential resets:
\begin{equation}
    S(t) = \begin{cases} 1 & \text{if } U(t) \geq \vartheta, \\ 0 & \text{otherwise}, \end{cases}
    \qquad
    U(t^+) = U_{\mathrm{reset}} \quad \text{if } S(t) = 1.
    \label{eq:lif-spike}
 from report.tex

\chapter{Background and Related Work}
\label{ch:background}

This chapter establishes the technical foundations on which this thesis is built. Section~\ref{sec:esc} introduces the environmental sound classification task and the ESC-50 benchmark. Section~\ref{sec:snn-background} formalises the spiking neural network framework, covering the \gls{lif} neuron model, surrogate gradient training, and the seven spike encoding methods evaluated in this work. Section~\ref{sec:neuromorphic-hardware} describes the SpiNNaker neuromorphic platform and situates it among contemporary alternatives. Section~\ref{sec:energy-background} examines the energy efficiency argument for \glspl{snn}, including the conditions under which spiking architectures genuinely outperform conventional networks. Section~\ref{sec:related-audio-snn} surveys prior work on \gls{snn}-based audio classification and identifies the gap this thesis addresses. Finally, Section~\ref{sec:adversarial-cl} reviews findings on adversarial robustness and continual learning in \glspl{snn}, motivating two of the advanced experiments in this report.

%=================================================================
\section{Environmental Sound Classification}
\label{sec:esc}
%=================================================================

Environmental sound classification (ESC) is a core task within computational auditory scene analysis, requiring a system to assign a categorical label to a short audio recording of a real-world acoustic event~\citep{barchiesi2015}. Applications include noise pollution monitoring in smart cities~\citep{bello2019}, automated wildlife census~\citep{stowell2019}, assistive technology for the hearing-impaired, and anomaly detection in industrial settings. Unlike speech recognition, ESC must handle the full diversity of non-linguistic acoustic events---from a dog barking to a chainsaw operating---without the benefit of structured phonological constraints.

\paragraph{The ESC-50 dataset.}
\gls{esc}~\citep{piczak2015} is the standard benchmark for this task. It comprises 2{,}000 five-second recordings drawn from the Freesound repository, distributed uniformly across 50 classes (40 clips per class). The classes span five broad super-categories: animals (e.g.\ dog, rooster, crow), natural soundscapes (e.g.\ rain, sea waves, thunderstorm), human non-speech sounds (e.g.\ clapping, coughing, laughing), domestic/indoor sounds (e.g.\ clock alarm, door knock, vacuum cleaner), and urban/outdoor sounds (e.g.\ helicopter, siren, car horn). Crucially, the dataset defines five predefined folds for cross-validation, ensuring that results are directly comparable across studies. Human classification performance on ESC-50 is 81.3\%~\citep{piczak2015}, reflecting the genuine perceptual difficulty of some distinctions (e.g.\ sea waves versus rain, insects versus crackling fire).

The current state of the art is 99.1\%, achieved by OmniVec2~\citep{srivastava2024} using a multi-modal Transformer pretrained on diverse data. Among audio-only methods, Pretrained Audio Neural Networks~(\gls{panns}) with CNN14 fine-tuning achieve 94.7\%~\citep{kong2020}, and Audio Spectrogram Transformers reach 95.6\%~\citep{gong2021}. These results rely on large-scale pretraining on AudioSet~\citep{gemmeke2017}, a dataset of approximately 2 million 10-second clips spanning 527 event categories. Without external pretraining, the best reported \gls{ann} performance on ESC-50 is substantially lower, motivating our choice to evaluate scratch-trained models as the primary comparison.

\paragraph{Input representation.}
The dominant input representation for ESC is the log-mel spectrogram, computed by applying the short-time Fourier transform to the audio waveform, mapping the resulting frequency bins through a bank of triangular mel-scale filters that approximate the frequency resolution of the human cochlea, and converting power to a logarithmic decibel scale. In this work, we extract spectrograms with 64 mel bands at a sampling rate of 22{,}050\,Hz and a hop length of 512 samples, yielding a representation of shape $(1, 64, 216)$ per clip---one channel, 64 frequency bins, and 216 time frames. Values are normalised to $[0, 1]$ per sample, a requirement for the spike encoding methods described in Section~\ref{subsec:encoding}.

\paragraph{Broader context.}
Other datasets complement ESC-50. UrbanSound8K~\citep{salamon2014} provides 8{,}732 clips across 10 urban sound classes; it is smaller in vocabulary and widely considered less challenging. AudioSet~\citep{gemmeke2017} provides massive scale (527 classes, 2 million clips) but uses weak labels and variable-quality recordings, making it unsuitable as a controlled evaluation benchmark. The Spiking Heidelberg Digits (SHD) and Spiking Speech Commands datasets~\citep{cramer2020} are natively spiking but address spoken-digit recognition rather than environmental sounds. No standard spiking benchmark exists for ESC, underscoring the need for this work.


%=================================================================
\section{Spiking Neural Networks}
\label{sec:snn-background}
%=================================================================

\Glspl{snn} are neural networks in which information is communicated through discrete, binary events called spikes. Unlike conventional \glspl{ann}, where neurons transmit continuous-valued activations, \gls{snn} neurons accumulate incoming signals over time in a membrane potential and emit a spike only when that potential exceeds a threshold. This event-driven, temporal computation paradigm is inspired by biological neural circuits and offers potential advantages in energy efficiency and temporal processing~\citep{maass1997,gerstner2002}.

%-------------------------------
\subsection{The Leaky Integrate-and-Fire Neuron}
\label{subsec:lif}
%-------------------------------

The \gls{lif} neuron is the most widely used model in modern \gls{snn} research, balancing biological plausibility with computational tractability~\citep{gerstner2002}. Its continuous-time dynamics describe an RC circuit with passive membrane leakage:

\begin{definition}[Leaky Integrate-and-Fire Neuron]
\label{def:lif}
The membrane potential $U(t)$ of a \gls{lif} neuron evolves according to the differential equation
\begin{equation}
    \tau_m \frac{dU(t)}{dt} = -U(t) + R \, I(t),
    \label{eq:lif-continuous}
\end{equation}
where $\tau_m$ is the membrane time constant, $R$ is the membrane resistance, and $I(t)$ is the input current. When $U(t)$ reaches the firing threshold $\vartheta$, the neuron emits a spike and the membrane potential resets:
\begin{equation}
    S(t) = \begin{cases} 1 & \text{if } U(t) \geq \vartheta, \\ 0 & \text{otherwise}, \end{cases}
    \qquad
    U(t^+) = U_{\mathrm{reset}} \quad \text{if } S(t) = 1.
    \label{eq:lif-spike}
 from report.tex

\chapter{Background and Related Work}
\label{ch:background}

This chapter establishes the technical foundations on which this thesis is built. Section~\ref{sec:esc} introduces the environmental sound classification task and the ESC-50 benchmark. Section~\ref{sec:snn-background} formalises the spiking neural network framework, covering the \gls{lif} neuron model, surrogate gradient training, and the seven spike encoding methods evaluated in this work. Section~\ref{sec:neuromorphic-hardware} describes the SpiNNaker neuromorphic platform and situates it among contemporary alternatives. Section~\ref{sec:energy-background} examines the energy efficiency argument for \glspl{snn}, including the conditions under which spiking architectures genuinely outperform conventional networks. Section~\ref{sec:related-audio-snn} surveys prior work on \gls{snn}-based audio classification and identifies the gap this thesis addresses. Finally, Section~\ref{sec:adversarial-cl} reviews findings on adversarial robustness and continual learning in \glspl{snn}, motivating two of the advanced experiments in this report.

%=================================================================
\section{Environmental Sound Classification}
\label{sec:esc}
%=================================================================

Environmental sound classification (ESC) is a core task within computational auditory scene analysis, requiring a system to assign a categorical label to a short audio recording of a real-world acoustic event~\citep{barchiesi2015}. Applications include noise pollution monitoring in smart cities~\citep{bello2019}, automated wildlife census~\citep{stowell2019}, assistive technology for the hearing-impaired, and anomaly detection in industrial settings. Unlike speech recognition, ESC must handle the full diversity of non-linguistic acoustic events---from a dog barking to a chainsaw operating---without the benefit of structured phonological constraints.

\paragraph{The ESC-50 dataset.}
\gls{esc}~\citep{piczak2015} is the standard benchmark for this task. It comprises 2{,}000 five-second recordings drawn from the Freesound repository, distributed uniformly across 50 classes (40 clips per class). The classes span five broad super-categories: animals (e.g.\ dog, rooster, crow), natural soundscapes (e.g.\ rain, sea waves, thunderstorm), human non-speech sounds (e.g.\ clapping, coughing, laughing), domestic/indoor sounds (e.g.\ clock alarm, door knock, vacuum cleaner), and urban/outdoor sounds (e.g.\ helicopter, siren, car horn). Crucially, the dataset defines five predefined folds for cross-validation, ensuring that results are directly comparable across studies. Human classification performance on ESC-50 is 81.3\%~\citep{piczak2015}, reflecting the genuine perceptual difficulty of some distinctions (e.g.\ sea waves versus rain, insects versus crackling fire).

The current state of the art is 99.1\%, achieved by OmniVec2~\citep{srivastava2024} using a multi-modal Transformer pretrained on diverse data. Among audio-only methods, Pretrained Audio Neural Networks~(\gls{panns}) with CNN14 fine-tuning achieve 94.7\%~\citep{kong2020}, and Audio Spectrogram Transformers reach 95.6\%~\citep{gong2021}. These results rely on large-scale pretraining on AudioSet~\citep{gemmeke2017}, a dataset of approximately 2 million 10-second clips spanning 527 event categories. Without external pretraining, the best reported \gls{ann} performance on ESC-50 is substantially lower, motivating our choice to evaluate scratch-trained models as the primary comparison.

\paragraph{Input representation.}
The dominant input representation for ESC is the log-mel spectrogram, computed by applying the short-time Fourier transform to the audio waveform, mapping the resulting frequency bins through a bank of triangular mel-scale filters that approximate the frequency resolution of the human cochlea, and converting power to a logarithmic decibel scale. In this work, we extract spectrograms with 64 mel bands at a sampling rate of 22{,}050\,Hz and a hop length of 512 samples, yielding a representation of shape $(1, 64, 216)$ per clip---one channel, 64 frequency bins, and 216 time frames. Values are normalised to $[0, 1]$ per sample, a requirement for the spike encoding methods described in Section~\ref{subsec:encoding}.

\paragraph{Broader context.}
Other datasets complement ESC-50. UrbanSound8K~\citep{salamon2014} provides 8{,}732 clips across 10 urban sound classes; it is smaller in vocabulary and widely considered less challenging. AudioSet~\citep{gemmeke2017} provides massive scale (527 classes, 2 million clips) but uses weak labels and variable-quality recordings, making it unsuitable as a controlled evaluation benchmark. The Spiking Heidelberg Digits (SHD) and Spiking Speech Commands datasets~\citep{cramer2020} are natively spiking but address spoken-digit recognition rather than environmental sounds. No standard spiking benchmark exists for ESC, underscoring the need for this work.


%=================================================================
\section{Spiking Neural Networks}
\label{sec:snn-background}
%=================================================================

\Glspl{snn} are neural networks in which information is communicated through discrete, binary events called spikes. Unlike conventional \glspl{ann}, where neurons transmit continuous-valued activations, \gls{snn} neurons accumulate incoming signals over time in a membrane potential and emit a spike only when that potential exceeds a threshold. This event-driven, temporal computation paradigm is inspired by biological neural circuits and offers potential advantages in energy efficiency and temporal processing~\citep{maass1997,gerstner2002}.

%-------------------------------
\subsection{The Leaky Integrate-and-Fire Neuron}
\label{subsec:lif}
%-------------------------------

The \gls{lif} neuron is the most widely used model in modern \gls{snn} research, balancing biological plausibility with computational tractability~\citep{gerstner2002}. Its continuous-time dynamics describe an RC circuit with passive membrane leakage:

\begin{definition}[Leaky Integrate-and-Fire Neuron]
\label{def:lif}
The membrane potential $U(t)$ of a \gls{lif} neuron evolves according to the differential equation
\begin{equation}
    \tau_m \frac{dU(t)}{dt} = -U(t) + R \, I(t),
    \label{eq:lif-continuous}
\end{equation}
where $\tau_m$ is the membrane time constant, $R$ is the membrane resistance, and $I(t)$ is the input current. When $U(t)$ reaches the firing threshold $\vartheta$, the neuron emits a spike and the membrane potential resets:
\begin{equation}
    S(t) = \begin{cases} 1 & \text{if } U(t) \geq \vartheta, \\ 0 & \text{otherwise}, \end{cases}
    \qquad
    U(t^+) = U_{\mathrm{reset}} \quad \text{if } S(t) = 1.
    \label{eq:lif-spike}
 from report.tex

\chapter{Background and Related Work}
\label{ch:background}

This chapter establishes the technical foundations on which this thesis is built. Section~\ref{sec:esc} introduces the environmental sound classification task and the ESC-50 benchmark. Section~\ref{sec:snn-background} formalises the spiking neural network framework, covering the \gls{lif} neuron model, surrogate gradient training, and the seven spike encoding methods evaluated in this work. Section~\ref{sec:neuromorphic-hardware} describes the SpiNNaker neuromorphic platform and situates it among contemporary alternatives. Section~\ref{sec:energy-background} examines the energy efficiency argument for \glspl{snn}, including the conditions under which spiking architectures genuinely outperform conventional networks. Section~\ref{sec:related-audio-snn} surveys prior work on \gls{snn}-based audio classification and identifies the gap this thesis addresses. Finally, Section~\ref{sec:adversarial-cl} reviews findings on adversarial robustness and continual learning in \glspl{snn}, motivating two of the advanced experiments in this report.

%=================================================================
\section{Environmental Sound Classification}
\label{sec:esc}
%=================================================================

Environmental sound classification (ESC) is a core task within computational auditory scene analysis, requiring a system to assign a categorical label to a short audio recording of a real-world acoustic event~\citep{barchiesi2015}. Applications include noise pollution monitoring in smart cities~\citep{bello2019}, automated wildlife census~\citep{stowell2019}, assistive technology for the hearing-impaired, and anomaly detection in industrial settings. Unlike speech recognition, ESC must handle the full diversity of non-linguistic acoustic events---from a dog barking to a chainsaw operating---without the benefit of structured phonological constraints.

\paragraph{The ESC-50 dataset.}
\gls{esc}~\citep{piczak2015} is the standard benchmark for this task. It comprises 2{,}000 five-second recordings drawn from the Freesound repository, distributed uniformly across 50 classes (40 clips per class). The classes span five broad super-categories: animals (e.g.\ dog, rooster, crow), natural soundscapes (e.g.\ rain, sea waves, thunderstorm), human non-speech sounds (e.g.\ clapping, coughing, laughing), domestic/indoor sounds (e.g.\ clock alarm, door knock, vacuum cleaner), and urban/outdoor sounds (e.g.\ helicopter, siren, car horn). Crucially, the dataset defines five predefined folds for cross-validation, ensuring that results are directly comparable across studies. Human classification performance on ESC-50 is 81.3\%~\citep{piczak2015}, reflecting the genuine perceptual difficulty of some distinctions (e.g.\ sea waves versus rain, insects versus crackling fire).

The current state of the art is 99.1\%, achieved by OmniVec2~\citep{srivastava2024} using a multi-modal Transformer pretrained on diverse data. Among audio-only methods, Pretrained Audio Neural Networks~(\gls{panns}) with CNN14 fine-tuning achieve 94.7\%~\citep{kong2020}, and Audio Spectrogram Transformers reach 95.6\%~\citep{gong2021}. These results rely on large-scale pretraining on AudioSet~\citep{gemmeke2017}, a dataset of approximately 2 million 10-second clips spanning 527 event categories. Without external pretraining, the best reported \gls{ann} performance on ESC-50 is substantially lower, motivating our choice to evaluate scratch-trained models as the primary comparison.

\paragraph{Input representation.}
The dominant input representation for ESC is the log-mel spectrogram, computed by applying the short-time Fourier transform to the audio waveform, mapping the resulting frequency bins through a bank of triangular mel-scale filters that approximate the frequency resolution of the human cochlea, and converting power to a logarithmic decibel scale. In this work, we extract spectrograms with 64 mel bands at a sampling rate of 22{,}050\,Hz and a hop length of 512 samples, yielding a representation of shape $(1, 64, 216)$ per clip---one channel, 64 frequency bins, and 216 time frames. Values are normalised to $[0, 1]$ per sample, a requirement for the spike encoding methods described in Section~\ref{subsec:encoding}.

\paragraph{Broader context.}
Other datasets complement ESC-50. UrbanSound8K~\citep{salamon2014} provides 8{,}732 clips across 10 urban sound classes; it is smaller in vocabulary and widely considered less challenging. AudioSet~\citep{gemmeke2017} provides massive scale (527 classes, 2 million clips) but uses weak labels and variable-quality recordings, making it unsuitable as a controlled evaluation benchmark. The Spiking Heidelberg Digits (SHD) and Spiking Speech Commands datasets~\citep{cramer2020} are natively spiking but address spoken-digit recognition rather than environmental sounds. No standard spiking benchmark exists for ESC, underscoring the need for this work.


%=================================================================
\section{Spiking Neural Networks}
\label{sec:snn-background}
%=================================================================

\Glspl{snn} are neural networks in which information is communicated through discrete, binary events called spikes. Unlike conventional \glspl{ann}, where neurons transmit continuous-valued activations, \gls{snn} neurons accumulate incoming signals over time in a membrane potential and emit a spike only when that potential exceeds a threshold. This event-driven, temporal computation paradigm is inspired by biological neural circuits and offers potential advantages in energy efficiency and temporal processing~\citep{maass1997,gerstner2002}.

%-------------------------------
\subsection{The Leaky Integrate-and-Fire Neuron}
\label{subsec:lif}
%-------------------------------

The \gls{lif} neuron is the most widely used model in modern \gls{snn} research, balancing biological plausibility with computational tractability~\citep{gerstner2002}. Its continuous-time dynamics describe an RC circuit with passive membrane leakage:

\begin{definition}[Leaky Integrate-and-Fire Neuron]
\label{def:lif}
The membrane potential $U(t)$ of a \gls{lif} neuron evolves according to the differential equation
\begin{equation}
    \tau_m \frac{dU(t)}{dt} = -U(t) + R \, I(t),
    \label{eq:lif-continuous}
\end{equation}
where $\tau_m$ is the membrane time constant, $R$ is the membrane resistance, and $I(t)$ is the input current. When $U(t)$ reaches the firing threshold $\vartheta$, the neuron emits a spike and the membrane potential resets:
\begin{equation}
    S(t) = \begin{cases} 1 & \text{if } U(t) \geq \vartheta, \\ 0 & \text{otherwise}, \end{cases}
    \qquad
    U(t^+) = U_{\mathrm{reset}} \quad \text{if } S(t) = 1.
    \label{eq:lif-spike}
